\documentclass[12pt]{article}
\usepackage[a4paper,margin=1.25cm]{geometry}
\usepackage{amsmath,amssymb,mathtools}
\usepackage{graphicx}
\usepackage{booktabs}
\usepackage{tabularx}
\usepackage{enumitem}
\usepackage{tikz}
\usepackage[T1]{fontenc}
\usepackage[utf8]{inputenc}
\usepackage[english]{babel}

\setlength{\parindent}{0pt}
\setlength{\parskip}{0.45em}
\setlength{\abovedisplayskip}{5pt plus 1pt minus 1pt}
\setlength{\belowdisplayskip}{5pt plus 1pt minus 1pt}
\setlength{\abovedisplayshortskip}{3pt plus 1pt minus 1pt}
\setlength{\belowdisplayshortskip}{3pt plus 1pt minus 1pt}

\newcommand{\dd}{\,\mathrm{d}}
\newcommand{\e}{\mathrm{e}}
\newcommand{\sumi}{\sum_{i=1}^{n}}
\newcommand{\W}{\operatorname{width}}
\newcommand{\fr}{\operatorname{frac}}
\newcommand{\dist}{\operatorname{dist}}
\newcommand{\ZZ}{\mathbb Z}
\newcommand{\normZ}[1]{\left\lVert #1\right\rVert_{\ZZ}}

\begin{document}

\begin{center}
{\Large\bfseries The width of an \(n\)-good array}\\[0.35em]
{\large golden conjugates, an unreachable integer, and how far you must miss it}
\end{center}

\textbf{Problem.} For an integer \(n\ge4\), call a real tuple
\((x_1,\dots,x_n)\) an \emph{\(n\)-good array} if
\[
\sumi x_i=n,\qquad \sumi x_i^2=2n,\qquad \sumi x_i^3=3n,
\]
and let \(\W(x)=\max_i x_i-\min_i x_i\).
\begin{enumerate}
\item[(1)] Find the largest real \(C\) with \(\W(x)\ge C\) for every
\(n\ge4\) and every \(n\)-good array.
\item[(2)] For that \(C\), prove that some \(\lambda>0\) satisfies
\(\W(x)>C+\lambda n^{-3/2}\) for every \(n\ge4\) and every \(n\)-good
array.
\end{enumerate}

\[
\boxed{\displaystyle
C=\sqrt5;
\qquad
\W(x)>\sqrt5+\frac{\lambda}{n^{3/2}}
\ \text{ holds with }\ \lambda=\tfrac18. }
\]
Sections 1--9 prove both statements. Sections 10--13 are a separate,
numerical study of how far the exponent \(\frac32\) and the constant
\(\lambda\) can be pushed; the results there are stated as conjectures,
with the evidence given.

\textbf{Overview.}
The three conditions are not independent constraints to be juggled: they
are a disguised statement about the golden ratio. Let
\[
\varphi=\frac{1+\sqrt5}{2},\qquad a=\frac{1-\sqrt5}{2},
\]
the two roots of \(t^2=t+1\), so that also \(t^3=2t+1\) for \(t\in\{a,\varphi\}\).
If \emph{every} \(x_i\) lies in \(\{a,\varphi\}\), then
\[
\sumi x_i^2=\sumi x_i+n,
\qquad
\sumi x_i^3=2\sumi x_i+n,
\]
so the first condition \(\sum x_i=n\) forces the other two automatically.
Such an array is \(n\)-good, has width \(\varphi-a=\sqrt5\), and is the
only possible way to achieve width \(\sqrt5\).

But it cannot exist. If \(k\) entries equal \(\varphi\) and \(n-k\) equal
\(a\), then \(\sum x_i=n\) reads \(k\sqrt5=n(1-a)\), i.e.
\[
k=np,\qquad p:=\frac{1-a}{\sqrt5}=\frac{5+\sqrt5}{10}=0.7236\ldots,
\]
and \(np\) is irrational, never an integer. The ideal configuration is
arithmetically forbidden, so \(\W>\sqrt5\) strictly. It can still be
approached, because \(np\) can be made arbitrarily close to an integer.
That is part (1).

Part (2) asks how badly the near-misses must fail, and the answer is
governed by how close \(np\) can be to \(\ZZ\). Since \(p\) is a
quadratic irrational, \(\normZ{np}\ge\frac2{5n}\); and an array of width
\(\sqrt5(1+\eta)\) can be shown to force \(\normZ{np}\ll n^2\eta^2\).
Comparing the two gives \(\eta\gg n^{-3/2}\).

Sections 10--13 then ask what the narrowest \(n\)-good array actually
looks like. Numerically it has three distinct values, and its width
follows a parametric law
\(M_n=\sqrt5+\sqrt5\,\Phi(\{np\})/n+O(n^{-2})\) rather than a single power
of \(n\): the excess is of order \(\frac1n\) for most \(n\) and drops to
order \(n^{-3/2}\) only along the odd-index Lucas numbers, where
\(\{np\}\) is closest to \(1\). None of this is proved here, and Section
12 states precisely what would be needed.

\section*{1. Normalisation: send \(\{a,\varphi\}\) to \(\{0,1\}\)}

Put
\[
x_i=a+\sqrt5\,u_i,
\qquad\text{equivalently}\qquad
u_i=\frac{x_i-a}{\sqrt5},
\qquad
\W(x)=\sqrt5\,\W(u),
\tag{1}
\]
so that \(x_i=a\) corresponds to \(u_i=0\) and \(x_i=\varphi\) to
\(u_i=1\).

\textbf{Lemma 1.} With \(p=\frac{5+\sqrt5}{10}\), the tuple \(x\) is
\(n\)-good if and only if
\[
\boxed{\displaystyle \sumi u_i=\sumi u_i^2=\sumi u_i^3=np .}
\tag{2}
\]

\textit{Proof.} Use \(a^2=a+1\) and \(a^3=2a+1\). Expanding (1),
\[
\sumi x_i=na+\sqrt5\sumi u_i,
\]
so \(\sum x_i=n\) is \(\sum u_i=\frac{n(1-a)}{\sqrt5}=np\). Next,
\[
\sumi x_i^2=na^2+2a\sqrt5\sumi u_i+5\sumi u_i^2 ,
\]
and substituting \(a^2=a+1\) and \(\sum u_i=np=\frac{n(1-a)}{\sqrt5}\),
\[
5\sumi u_i^2=2n-n(a+1)-2an(1-a)=n(1-a)(1-2a)=n(1-a)\sqrt5,
\]
since \(1-2a=\sqrt5\); hence \(\sum u_i^2=\frac{n(1-a)}{\sqrt5}=np\).
Similarly
\[
\sumi x_i^3=na^3+3a^2\sqrt5\sumi u_i+15a\sumi u_i^2+5\sqrt5\sumi u_i^3 ,
\]
and using \(a^3=2a+1\), \(a^2=a+1\) and \(3a(1+\sqrt5)=-6\),
\[
5\sqrt5\sumi u_i^3=n(1-a)\bigl[2-3a^2-3\sqrt5\,a\bigr]=5n(1-a),
\]
so \(\sum u_i^3=np\). Every step is reversible. \(\square\)

The three inhomogeneous conditions have become one homogeneous statement:
\emph{the first three power sums of \(u\) coincide}. The specific value
\(np\) they share is the entire arithmetic content of the problem.

\section*{2. The test polynomial}

Let
\[
q(t)=t(1-t).
\]
By (2),
\[
\sumi q(u_i)=\sumi u_i-\sumi u_i^2=0,
\qquad
\sumi u_i\,q(u_i)=\sumi u_i^2-\sumi u_i^3=0 .
\tag{3}
\]
So \(q\) has mean zero against both the counting measure and the measure
weighted by \(u\). Since \(q\ge0\) exactly on \([0,1]\), the first
identity in (3) already says that the \(u_i\) cannot all lie strictly
inside \([0,1]\) unless they sit at the endpoints. The two identities
together pin both ends.

\textbf{Lemma 2.} Let \(m=\min_i u_i\) and \(M=\max_i u_i\). Then
\(M\ge1\) and \(m\le0\).

\textit{Proof.} Suppose \(M<1\). Then \(1-u_i>0\) for all \(i\), so by (3)
\[
0=\sumi u_i\,q(u_i)=\sumi u_i^2(1-u_i)
\]
is a sum of non-negative terms; hence \(u_i^2(1-u_i)=0\), i.e. \(u_i=0\)
for every \(i\). But then \(\sum u_i=0\ne np\).

Suppose \(m>0\). Subtracting the two identities in (3),
\[
0=\sumi(u_i-1)q(u_i)=-\sumi u_i(u_i-1)^2 ,
\]
again a sum of non-negative terms since \(u_i>0\); hence \(u_i=1\) for
every \(i\), and \(\sum u_i=n\ne np\) because \(p\ne1\). \(\square\)

\section*{3. Part (1), lower bound: \(\W(x)>\sqrt5\)}

By Lemma 2,
\[
\W(u)=M-m\ge1-0=1 .
\]
Suppose \(\W(u)=1\). Then \(m=0\) and \(M=1\) (Lemma 2 forces
\(m\le0\le1\le M\), and \(M-m=1\) leaves no slack), so
\(0\le u_i\le1\) for all \(i\). On \([0,1]\) we have \(q\ge0\), and
\(\sum q(u_i)=0\) by (3), so \(q(u_i)=0\) for every \(i\), i.e.
\[
u_i\in\{0,1\}\quad\text{for all }i .
\]
Then \(\sum u_i\) is a non-negative integer. But (2) gives
\[
\sumi u_i=np=n\cdot\frac{5+\sqrt5}{10}\notin\ZZ,
\]
since \(\sqrt5\) is irrational and \(n\ge1\). Contradiction. Therefore
\[
\boxed{\displaystyle \W(u)>1,
\qquad
\W(x)=\sqrt5\,\W(u)>\sqrt5 .}
\tag{4}
\]

Unwinding the normalisation, the excluded configuration is exactly the
one described in the overview: all entries in \(\{a,\varphi\}\), which
would be \(n\)-good but requires \(np\in\ZZ\).

\section*{4. Part (1), sharpness: a three-point patch}

It remains to show that \(\sqrt5\) is approached, so that no larger
constant works. The idea is to keep most entries at \(0\) and \(1\) and
absorb the fractional part of \(np\) into three extra entries.

\textbf{Lemma 3.} For every \(R\in(1,2)\) there exist three distinct
reals \(y_1,y_2,y_3\) with
\[
y_1+y_2+y_3=y_1^2+y_2^2+y_3^2=y_1^3+y_2^3+y_3^3=R .
\tag{5}
\]

\textit{Proof.} Prescribe the power sums \(s_1=s_2=s_3=R\) and read off
the elementary symmetric functions from Newton's identities:
\[
e_1=R,\qquad
e_2=\frac{s_1e_1-s_2}{2}=\frac{R^2-R}{2},\qquad
e_3=\frac{e_2s_1-e_1s_2+s_3}{3}=\frac{R(R-1)(R-2)}{6}.
\]
The corresponding monic cubic is
\[
F_R(y)=y^3-Ry^2+\frac{R^2-R}{2}\,y-\frac{R(R-1)(R-2)}{6},
\tag{6}
\]
whose discriminant is
\[
\Delta(F_R)=\frac{R^2(3-R)^2(2-R)(R-1)}{6} .
\tag{7}
\]
For \(1<R<2\) all four factors are positive, so \(\Delta>0\) and \(F_R\)
has three distinct real roots; by construction they satisfy (5).
\(\square\)

(Identity (7) is the ordinary cubic discriminant
\(18e_1e_2e_3-4e_1^3e_3+e_1^2e_2^2-4e_2^3-27e_3^2\) after simplification;
it was verified symbolically and numerically.)

\textbf{The construction.} Since \(p\) is irrational, \(\{np\}_{n\ge1}\)
is equidistributed \(\bmod\,1\); more concretely, if \(p_k/q_k\) are the
continued-fraction convergents of \(p\), then for those \(k\) with
\(q_kp-p_k>0\) one has
\[
0<\fr(q_kp)=q_kp-p_k<\frac1{q_k},
\]
and these occur for infinitely many \(k\). So choose \(n_j\to\infty\)
with
\[
r_j:=\fr(n_jp)\longrightarrow0^+ ,
\]
and set
\[
K_j=\lfloor n_jp\rfloor-1,\qquad R_j=1+r_j\in(1,2),
\qquad\text{so that}\qquad K_j+R_j=n_jp .
\tag{8}
\]
For large \(j\) both \(K_j\ge0\) and \(n_j-K_j-3\ge0\) (as
\(K_j\approx0.7236\,n_j\)). Take the array
\[
u^{(j)}=\bigl(\underbrace{1,\dots,1}_{K_j},\
\underbrace{0,\dots,0}_{n_j-K_j-3},\ y_1,y_2,y_3\bigr),
\qquad y_i \text{ the roots of }F_{R_j}.
\]
Then for \(k=1,2,3\),
\[
\sum_i \bigl(u^{(j)}_i\bigr)^k=K_j\cdot1+0+R_j=n_jp,
\]
so by Lemma 1 the corresponding \(x^{(j)}\) is \(n_j\)-good.

As \(R_j\to1^+\), the cubic (6) tends to \(y^3-y^2=y^2(y-1)\), so the root
multiset converges: \(\{y_1,y_2,y_3\}\to\{0,0,1\}\). Hence
\(\W(u^{(j)})\to1\) and
\[
\W(x^{(j)})=\sqrt5\,\W(u^{(j)})\longrightarrow\sqrt5 .
\]
Combined with (4),
\[
\boxed{\displaystyle
\inf\W(x)=\sqrt5,
\qquad
C=\sqrt5 . }
\]

This settles part (1). The patched array is not the narrowest \(n\)-good
array, though: it spends three entries on the cubic, and Section 10
exhibits a minimiser using only one interior value, whose excess over
\(\sqrt5\) is smaller by a factor of order \(n\).

\begin{center}
\small
\renewcommand{\arraystretch}{1.35}
\begin{tabularx}{0.98\textwidth}{@{}rXXXX@{}}
\toprule
\(n\) & \(\fr(np)\) & \(\W(u)\) & \(\eta=\W(u)-1\) & \(\W(x)\) \\
\midrule
\(14\)      & \(1.30\cdot10^{-1}\) & \(1.1448171\) & \(1.45\cdot10^{-1}\) & \(2.5598889\) \\
\(123\)     & \(3.64\cdot10^{-3}\) & \(1.0246280\) & \(2.46\cdot10^{-2}\) & \(2.2911378\) \\
\(1165\)    & \(1.92\cdot10^{-3}\) & \(1.0178901\) & \(1.79\cdot10^{-2}\) & \(2.2760715\) \\
\(15127\)   & \(2.96\cdot10^{-5}\) & \(1.0022198\) & \(2.22\cdot10^{-3}\) & \(2.2410315\) \\
\(103682\)  & \(4.31\cdot10^{-6}\) & \(1.0008479\) & \(8.48\cdot10^{-4}\) & \(2.2379639\) \\
\midrule
\(\infty\) & \(0\) & \(1\) & \(0\) & \(\sqrt5=2.2360680\) \\
\bottomrule
\end{tabularx}
\end{center}

The three power sums of each row were verified to \(25\) decimal places.

\section*{5. Part (2), setup}

Write
\[
\W(u)=M-m=1+\eta,
\qquad \eta>0 \text{ by } (4).
\]
If \(\eta\ge1\) then \(\W(x)\ge2\sqrt5>\sqrt5+\lambda n^{-3/2}\) for any
\(\lambda\le1\), and there is nothing to prove; so assume
\(0<\eta<1\) from now on.

By Lemma 2, \(m\le0\) and \(M\ge1\); together with \(M-m=1+\eta\) this
gives \(m\ge M-1-\eta\ge-\eta\) and \(M\le m+1+\eta\le1+\eta\), so
\[
\boxed{\displaystyle -\eta\le u_i\le1+\eta\quad\text{for all }i.}
\tag{9}
\]
Every entry is trapped in a window of width \(1+2\eta\) around
\([0,1]\). Write \(\normZ{t}=\min_{K\in\ZZ}|t-K|\) and
\(d_i=\dist(u_i,\{0,1\})\).

Sections 6 to 8 bound \(\sum d_i\) by \(n\eta\), upgrade that to a bound
on \(\sum d_i^2\), and convert it into an upper bound for \(\normZ{np}\).
Section 8 contrasts this with a lower bound coming from the irrationality
of \(\sqrt5\).

\section*{6. The entries are close to \(\{0,1\}\) in \(\ell^1\)}

On \([0,1]\), \(q(t)=t(1-t)=\min(t,1-t)\cdot\max(t,1-t)\ge\frac12 d(t)\),
where \(d(t)=\dist(t,\{0,1\})\); that is,
\[
d(t)\le2q(t),\qquad 0\le t\le1 .
\tag{10}
\]
On the two flanks \(t\in[-\eta,0]\cup[1,1+\eta]\),
\[
q(t)\ge-\eta(1+\eta)\ge-2\eta,
\qquad
d(t)\le\eta .
\tag{11}
\]
By (3) and (11), splitting according to whether \(u_i\in[0,1]\),
\[
\sum_{u_i\in[0,1]}q(u_i)
=-\sum_{u_i\notin[0,1]}q(u_i)
\le 2n\eta ,
\]
and therefore, by (10) and (11) again,
\[
\sumi d_i
=\sum_{u_i\in[0,1]}d_i+\sum_{u_i\notin[0,1]}d_i
\le 2\sum_{u_i\in[0,1]}q(u_i)+n\eta
\le5n\eta .
\tag{12}
\]
Since the \(d_i\) are non-negative,
\((\sum d_i)^2=\sum d_i^2+\text{(cross terms)}\ge\sum d_i^2\), so
\[
\sumi d_i^2\le\Bigl(\sumi d_i\Bigr)^2\le25\,n^2\eta^2 .
\tag{13}
\]

\section*{7. A doubly tangent cubic converts \(\ell^2\) into an integer}

The bound (13) is quadratic in the distances, so it must be paired with a
quantity that is itself quadratically small near \(\{0,1\}\). Take
\[
B(t)=3t^2-2t^3 .
\]
It satisfies \(B(0)=0\), \(B(1)=1\), and \(B'(t)=6t(1-t)\) vanishes at
both endpoints; equivalently
\[
B(t)=t^2(3-2t),
\qquad
1-B(t)=(1-t)^2(1+2t),
\]
which exhibits the double contact directly. Consequently, for
\(-1\le t\le2\),
\[
\dist\bigl(B(t),\{0,1\}\bigr)\le5\,\dist(t,\{0,1\})^2 ,
\tag{14}
\]
since \(|B(t)|\le5t^2\) for \(t\le\frac12\) and
\(|1-B(t)|\le5(1-t)^2\) for \(t\ge\frac12\) on that range.

The point of choosing \(B\) among all doubly tangent cubics is that its
sum is computable from (2):
\[
\sumi B(u_i)=3\sumi u_i^2-2\sumi u_i^3=3np-2np=np .
\tag{15}
\]

Now let \(b_i\in\{0,1\}\) be a nearest Boolean value to \(B(u_i)\), and
put \(K=\sum_i b_i\in\ZZ\). By (9) we have \(u_i\in[-1,2]\), so (14)
applies, and with (15), (13),
\[
\normZ{np}\le\Bigl|np-K\Bigr|
=\Bigl|\sumi B(u_i)-\sumi b_i\Bigr|
\le\sumi\dist\bigl(B(u_i),\{0,1\}\bigr)
\le5\sumi d_i^2
\le125\,n^2\eta^2 .
\tag{16}
\]

\section*{8. \(np\) cannot be too close to an integer}

\textbf{Lemma 4.} For every integer \(n\ge4\),
\[
\normZ{np}\ge\frac{2}{5\,n},
\qquad p=\frac{5+\sqrt5}{10}.
\tag{17}
\]

\textit{Proof.} Let \(K\in\ZZ\) and put \(m=2K-n\in\ZZ\), so that
\(m\equiv n\pmod 2\) and
\[
10\,|np-K|=|5n+n\sqrt5-10K|=|n\sqrt5-5m| .
\]
If \(|n\sqrt5-5m|\ge1\) then \(|np-K|\ge\frac1{10}\ge\frac2{5n}\) for
\(n\ge4\). Otherwise \(|n\sqrt5-5m|<1\), which forces \(5m>n\sqrt5-1>0\)
and
\[
n\sqrt5+5m<2\sqrt5\,n+1\le5n\qquad(n\ge2).
\]
Now use the parity of \(m\). If \(n\) and \(m\) are both even then
\(4\mid n^2\) and \(4\mid m^2\); if both are odd then \(n^2\equiv m^2\equiv1
\pmod 8\), so \(n^2-5m^2\equiv-4\equiv0\pmod4\). In either case
\(4\mid n^2-5m^2\), and \(n^2-5m^2\ne0\) because \(\sqrt5\) is
irrational. Hence \(|n^2-5m^2|\ge4\) and
\[
|n\sqrt5-5m|=\frac{5\,|n^2-5m^2|}{n\sqrt5+5m}\ge\frac{5\cdot4}{5n}=\frac4n ,
\]
so \(|np-K|\ge\frac{2}{5n}\). Taking the minimum over \(K\) gives (17).
\(\square\)

The divisibility \(4\mid n^2-5m^2\) is what upgrades the naive
\(|n^2-5m^2|\ge1\) to \(\ge4\), and it costs nothing: it is only the
observation that \(m\) and \(n\) have the same parity by construction.

Bound (17) is close to best possible. Numerically
\(\min_{4\le n\le3\cdot10^5}n\normZ{np}=0.4222912\), attained at
\(n=4\), against the constant \(\frac25=0.4\) of (17); Section 12 shows
that the true asymptotic constant is \(\frac1{\sqrt5}=0.4472136\).
Note that \(\normZ{np}\asymp n^{-1}\) is \emph{false} as a statement about
all \(n\): for most \(n\) the quantity \(\normZ{np}\) is of order one, and
the content of (17) is the lower bound alone.

\section*{9. Part (2), conclusion}

Combining (16) and (17), for \(0<\eta<1\) and \(n\ge4\),
\[
\frac{2}{5n}\le\normZ{np}\le125\,n^2\eta^2
\qquad\Longrightarrow\qquad
\eta^2\ge\frac{2}{625\,n^3},
\]
that is,
\[
\boxed{\displaystyle \eta\ge\frac{\sqrt2}{25\,n^{3/2}}
=\frac{0.0565685\ldots}{n^{3/2}} .}
\tag{18}
\]
(The case \(\eta\ge1\) was settled in Section 5, where the bound is
trivial.) Hence
\[
\W(x)=\sqrt5\,\W(u)=\sqrt5(1+\eta)
\ge\sqrt5+\frac{\sqrt{10}}{25\,n^{3/2}}
=\sqrt5+\frac{0.1264911\ldots}{n^{3/2}},
\]
and therefore, with the explicit choice \(\lambda=\frac18=0.125\),
\[
\boxed{\displaystyle
\W(x)>\sqrt5+\frac{\lambda}{n^{3/2}}
\qquad\text{for every }n\ge4\text{ and every }n\text{-good array}. }
\]

The chain (12)--(18) is quantitative at every step, so \(\lambda\) comes
out explicitly rather than from a compactness or contradiction
argument.



\section*{10. A candidate minimiser}

Sections 10--13 are not part of the proof. They describe the narrowest
\(n\)-good array that numerical search finds and work out its
asymptotics. The formulas are supported by computation and by a
consistent asymptotic expansion, but global optimality is not proved;
Section 12 states which steps are missing.

Write
\[
M_n=\inf\bigl\{\W(x):x\ \text{is }n\text{-good}\bigr\},
\qquad
W_n=\frac{M_n}{\sqrt5}=\inf\W(u),
\]
so that part (2) asserts \(M_n>\sqrt5+\lambda n^{-3/2}\).

The configuration found by numerical minimisation takes three distinct
values, with the two extreme ones carrying almost all the mass.

\textbf{The ansatz.} Fix \(n\) and put
\[
k=\lfloor np\rfloor,\qquad \ell=n-k-1,\qquad r=\{np\},
\]
and consider arrays of the form
\[
u=\bigl(\underbrace{v_-,\dots,v_-}_{\ell},\ w,\ \underbrace{v_+,\dots,v_+}_{k}\bigr),
\qquad v_-<w<v_+ .
\tag{19}
\]
The three moment equations (2) become three equations in \(v_-,w,v_+\),
\[
\ell v_-^{\,j}+w^{\,j}+k\,v_+^{\,j}=np,\qquad j=1,2,3,
\tag{20}
\]
so no free parameter remains. By Lemma 2 the extreme values satisfy
\(v_-\le0\) and \(v_+\ge1\), and by Section 3 at least one inequality is
strict; numerically both are.

\textbf{Reduction to one scalar equation.} System (20) is easier to
handle after eliminating two unknowns. For fixed \(w\), the first two
equations determine \(v_\pm\) uniquely: with \(T=n-1\) and \(s=np\),
\[
\mu(w)=\frac{s-w}{T},
\qquad
D(w)=s-w^2-\frac{(s-w)^2}{T},
\qquad
\delta(w)=\sqrt{\frac{T\,D(w)}{k\ell}},
\]
\[
v_-(w)=\mu(w)-\frac kT\delta(w),
\qquad
v_+(w)=\mu(w)+\frac \ell T\delta(w),
\tag{21}
\]
and the third equation becomes a single scalar condition
\[
G_n(w)=\ell\,v_-(w)^3+w^3+k\,v_+(w)^3-s=0 .
\tag{22}
\]
For every \(n\) checked, \(4\le n\le10^4\), one finds \(G_n(0)>0\),
\(G_n(1)<0\), and exactly one root in between, giving an ordered solution
\(v_-<0<w<1<v_+\). Proving these two sign conditions and the uniqueness
of the root would establish existence and uniqueness of (19); that is not
done here.

\textbf{Asymptotics.} Write \(v_-=-A/n\) and \(v_+=1+C/n\) and expand
(20) in \(1/n\), using \(\ell=(1-p)n+O(1)\) and \(k=pn+O(1)\). To leading
order,
\[
p\,C=w^2(1-w)+O(n^{-1}),
\qquad
(1-p)\,A=w(1-w)^2+O(n^{-1}),
\tag{23}
\]
and substituting back leaves a single condition on \(w\):
\[
\boxed{\displaystyle 3w^2-2w^3=\{np\}+O(n^{-1}).}
\tag{24}
\]
The error term in (24) is not negligible at small \(n\). At \(n=4\) the
actual root is \(w_4=0.8155117\ldots\), whereas the solution of
\(3w^2-2w^3=\{4p\}\) is \(0.7985\ldots\); the two agree only to first
order.

The cubic on the left of (24) is exactly the doubly tangent \(B\) of
Section 7, and the coincidence has a reason. Identity (15) says \(\sum_iB(u_i)=np\) for \emph{every} good array; evaluated on (19),
\(B(v_-)=O(n^{-2})\) and \(1-B(v_+)=O(n^{-2})\), so
\[
\ell B(v_-)+B(w)+kB(v_+)=B(w)+k+O(n^{-1})=np
\quad\Longrightarrow\quad
B(w)=\{np\}+O(n^{-1}).
\]
The function used to prove the lower bound is the same one that locates
the middle atom. Section 12 shows that this is why the bound of Section 9
is sharp in order.

\textbf{Inverting \(B\).} On \([0,1]\) the map \(B\) is an increasing
bijection onto \([0,1]\), and its inverse has a closed form:
\[
\boxed{\displaystyle
\sigma(r)=\frac12-\sin\!\left(\frac13\arcsin(1-2r)\right),
\qquad
B(\sigma(r))=r,\quad r\in[0,1].}
\tag{25}
\]
(Endpoints: \(r=0\) gives \(\frac12-\sin\frac\pi6=0\), \(r=1\) gives
\(\frac12+\sin\frac\pi6=1\), and \(r=\frac12\) gives \(\frac12\).)
Formula (25) is the trigonometric solution of \(2\sigma^3-3\sigma^2+r=0\),
which has three real roots for \(r\in(0,1)\); (25) selects the one in
\([0,1]\).

\section*{11. The parametric law}

Substituting (23) and (24) into \(\W(u)=v_+-v_-=1+\frac{A+C}{n}\) gives
the following description, again asymptotic and again unproved.

\textbf{Conjecture.} With \(\sigma=\sigma(\{np\})\) from (25) and
\[
\boxed{\displaystyle
\Phi(r)=\sigma(1-\sigma)\left(\frac{\sigma}{p}+\frac{1-\sigma}{1-p}\right),
\qquad 1-p=\frac{5-\sqrt5}{10},}
\tag{26}
\]
the minimal width satisfies
\[
\boxed{\displaystyle
W_n=1+\frac{\Phi(\{np\})}{n}+O\!\left(\frac1{n^{2}}\right),
\qquad
M_n=\sqrt5+\frac{\sqrt5\,\Phi(\{np\})}{n}+O\!\left(\frac1{n^{2}}\right).}
\tag{27}
\]

Two caveats on (27). It presumes that (19) is optimal, which is not
proved. And the remainder \(O(n^{-2})\) is claimed uniformly, whereas the
expansion leading to it is not uniform: \(B'(\sigma)\to0\) as
\(\{np\}\to0\) or \(1\), so the implicit-function step degenerates
precisely on the subsequences that matter in Section 12. Numerically the
remainder does behave like \(n^{-2}\) there, but the derivation above
does not establish it.

Granting (27), the answer is not a single power of \(n\): it is a fixed
profile \(\Phi\) sampled at the equidistributed points \(\{np\}\) and
scaled by \(\frac1n\). For most \(n\) the excess is of order \(\frac1n\);
smaller powers appear only when \(\{np\}\) approaches an endpoint.

Some values of \(\Phi\). Since
\[
p(1-p)=\frac{5+\sqrt5}{10}\cdot\frac{5-\sqrt5}{10}=\frac{20}{100}=\frac15,
\]
the midpoint value is exact:
\[
\Phi\!\left(\tfrac12\right)
=\frac14\cdot\frac12\left(\frac1p+\frac1{1-p}\right)
=\frac18\cdot\frac{1}{p(1-p)}
=\frac58 .
\]
The maximum is \(\Phi\approx0.65367\) near \(r\approx0.354\). Combined
with (27) this gives
\[
M_n-\sqrt5\le\frac{0.654\sqrt5}{n}+O(n^{-2}),
\]
not a clean bound for every \(n\): the remainder in (27) may be positive,
so the honest statement is that for each \(\varepsilon>0\) the excess is
at most \((0.654+\varepsilon)\sqrt5/n\) once \(n\) is large.

\textbf{The two edges.} Near \(r=0\) we have \(B(\sigma)\approx3\sigma^2\),
so \(\sigma\approx\sqrt{r/3}\) and only the second term of (26) survives:
\[
\Phi(r)=\frac1{1-p}\sqrt{\frac r3}+O(r),
\qquad r\to0^{+}.
\tag{28}
\]
Near \(r=1\), writing \(\tau=1-r\), the same expansion at the other end
gives \(1-\sigma\approx\sqrt{\tau/3}\) and
\[
\Phi(1-\tau)=\frac1p\sqrt{\frac\tau3}+O(\tau),
\qquad \tau\to0^{+}.
\tag{29}
\]
The two edges are not symmetric. Since \(p>1-p\), the coefficient
\(\frac1p\) in (29) is smaller than \(\frac1{1-p}\) in (28), so the
\emph{right} edge, \(\{np\}\) just below \(1\), produces the smaller
width. That is the dangerous side.

Combining (27) and (29), when \(np\) lies just below an integer,
\[
W_n-1
\;\sim\;
\frac{1}{p\sqrt3}\cdot\frac{\sqrt{\normZ{np}}}{n},
\tag{30}
\]
which is the square-root law behind the exponent \(\frac32\): a defect of
size \(\normZ{np}\) in the arithmetic costs only \(\sqrt{\normZ{np}}\) in
the geometry, because \(B\) is tangent to its endpoint values.

\section*{12. What is and is not established about the exponent}

Formula (30) turns the question into one about how well \(np\) can be
approximated by integers, and that part \emph{can} be settled exactly.

\textbf{An exact Diophantine identity.} Let \(\varphi=\frac{1+\sqrt5}2\),
\(\psi=\frac{1-\sqrt5}2=-\varphi^{-1}\), and let \(L_m=\varphi^m+\psi^m\),
\(F_m=\frac{\varphi^m-\psi^m}{\sqrt5}\) be the Lucas and Fibonacci
numbers. Since \(p=\frac12+\frac1{2\sqrt5}=\frac{\varphi}{\sqrt5}\),
\[
L_mp-F_{m+1}
=\frac{(\varphi^m+\psi^m)\varphi-(\varphi^{m+1}-\psi^{m+1})}{\sqrt5}
=\frac{\psi^m(\varphi+\psi)}{\sqrt5}
=\frac{(-1)^m}{\sqrt5\,\varphi^{m}},
\tag{31}
\]
using \(\varphi+\psi=1\). For \(m\ge1\) the right-hand side has modulus
below \(\frac12\), so \(F_{m+1}\) is the nearest integer and
\[
L_m\normZ{L_mp}=\frac{L_m}{\sqrt5\,\varphi^m}
=\frac{1+(-1)^m\varphi^{-2m}}{\sqrt5}.
\tag{32}
\]
For odd \(m\) this is \emph{below} \(\frac1{\sqrt5}\) and increases to it;
these are the \(n=L_m\) with \(\{np\}\to1^-\), the dangerous edge of
Section 11. The odd-index Lucas numbers are
\(4,11,29,76,199,521,1364,3571,9349,\dots\), satisfying
\(n_{j+1}=3n_j-n_{j-1}\). Together with Hurwitz's theorem, (32) gives
\[
\liminf_{n\to\infty}n\normZ{np}=\frac1{\sqrt5},
\tag{33}
\]
and it shows that the constant \(\frac25\) of Lemma 4 cannot be raised
above \(\frac1{\sqrt5}\).

\textbf{The conjectural constant.} Substituting
\(\normZ{np}\sim\frac1{\sqrt5\,n}\) into (30) and granting (27),
\[
\liminf_{n\to\infty}n^{3/2}\bigl(W_n-1\bigr)
\;\overset{?}{=}\;\frac1{p\sqrt{3\sqrt5}}=0.5335734771\ldots,
\tag{34}
\]
\[
\liminf_{n\to\infty}n^{3/2}\bigl(M_n-\sqrt5\bigr)
\;\overset{?}{=}\;\frac{\sqrt5}{p\sqrt{3\sqrt5}}=1.1931065657\ldots
\]
The question marks record that (34) inherits both unproved ingredients
of Section 11.

\textbf{Three statements that must be kept apart.} Let
\[
\Lambda_{\mathrm{unif}}=\inf_{n\ge4}n^{3/2}\bigl(M_n-\sqrt5\bigr)
\]
be the supremum of admissible constants in part (2). Then always
\[
\Lambda_{\mathrm{unif}}\le\liminf_{n\to\infty}n^{3/2}\bigl(M_n-\sqrt5\bigr),
\]
and equality is a separate claim: some finite \(n\), or some other family
of configurations, could give a smaller value. So even a proof of (34)
would not by itself identify the best \(\lambda\); it would only exhibit
an asymptotic obstruction. What Section 9 does prove is
\(\Lambda_{\mathrm{unif}}\ge\frac{\sqrt{10}}{25}=0.1264911\ldots\), a
factor \(9.43\) below the conjectured value.

\textbf{What would suffice for optimality of the exponent.} To show that
\(\frac32\) cannot be lowered, global optimality of (19) is \emph{not}
needed. It is enough to
\begin{enumerate}
\item[(a)] prove that (22) has a root \(w\in(0,1)\) with
\(v_-(w)<0<w<1<v_+(w)\), at least for \(n=L_m\) with \(m\) odd;
\item[(b)] prove for that explicit admissible array the bound
\(\W(x)-\sqrt5=O(n^{-3/2})\), using (31) for the size of \(\{np\}\).
\end{enumerate}
Then \(M_n-\sqrt5=O(n^{-3/2})\) along a subsequence, so no universal
bound \(\W(x)\ge\sqrt5+\lambda n^{-\alpha}\) with \(\alpha<\frac32\) can
hold, and part (2) asks for the correct power. Step (a) is a statement
about the sign of one explicit scalar function, which is a far smaller
task than a global minimisation theorem. Identifying \(\Lambda_{\rm
unif}\) or proving (27) would additionally require a genuine
moment-space argument: a dual polynomial certificate, or a complete
analysis of the optimality conditions ruling out configurations with two
interior atoms.

\textbf{Why the proof of Sections 6--9 is sharp in order.} Evaluate the
chain on the array (19) in the dangerous regime, where \(\sigma\to1\) and
\(\eta=W_n-1\sim(1-\sigma)/(pn)\).

\begin{center}
\small
\renewcommand{\arraystretch}{1.45}
\begin{tabularx}{\textwidth}{@{}>{\raggedright\arraybackslash}p{0.17\textwidth}XXX@{}}
\toprule
Quantity & proved bound & value on (19) & loss \\
\midrule
\(\sum_i d_i\)
& \(\le5n\eta\) \ (12)
& \(\sim2(1-\sigma)=2pn\eta\)
& \(3.5\times\) \\
\(\sum_i d_i^2\)
& \(\le(\sum_i d_i)^2\) \ (13)
& \(\sim(1-\sigma)^2=p^2n^2\eta^2\)
& \(4\times\) \\
\(\normZ{np}\)
& \(\le125n^2\eta^2\) \ (16)
& \(\sim3(1-\sigma)^2=3p^2n^2\eta^2\)
& \(79.6\times\) \\
\(\normZ{np}\)
& \(\ge\frac2{5n}\) \ (17)
& \(\sim\frac1{\sqrt5\,n}\) along odd Lucas \(n\)
& \(1.12\times\) \\
\bottomrule
\end{tabularx}
\end{center}

Every entry is tight in order of magnitude; only constants are lost, and
after Lemma 4 the arithmetic input is nearly sharp. The exact relation
behind the third row is
\[
1-B(\sigma)=(1-\sigma)^2(1+2\sigma),
\]
which is asymptotically \(3(1-\sigma)^2\) as \(\sigma\to1\); this is (24)
again, and the double tangency of \(B\) is the reason the exponent is
\(\frac32\) rather than \(2\). Feeding the sharp constants through the
same argument returns
\[
\eta\gtrsim\frac{1}{p\sqrt{3\sqrt5}}\,n^{-3/2},
\]
which is (34). The residual gap between the proved
\(\frac{\sqrt{10}}{25}\) and the conjectured
\(\frac{\sqrt5}{p\sqrt{3\sqrt5}}\) is a factor \(9.43\), splitting as
\(\sqrt{79.6}=8.92\) from the constants in (12), (13), (16) and
\(\sqrt{1.12}=1.06\) from Lemma 4.

\section*{13. Numerics}

Solving (20) to \(30\) digits gives the following. The dangerous
subsequence, \(n=L_m\) with \(m\) odd:

\begin{center}
\small
\renewcommand{\arraystretch}{1.3}
\begin{tabularx}{\textwidth}{@{}rXXXXX@{}}
\toprule
\(n\) & \(\{np\}\) & \(w_n\) & \(\W(u)-1\) & \(n^{3/2}(\W(u)-1)\) & \(n^{3/2}(\W(x)-\sqrt5)\) \\
\midrule
\(4\)    & \(0.894427\) & \(0.815512\) & \(8.7707\cdot10^{-2}\) & \(0.701658\) & \(1.568954\) \\
\(11\)   & \(0.959675\) & \(0.884207\) & \(1.6911\cdot10^{-2}\) & \(0.616974\) & \(1.379596\) \\
\(29\)   & \(0.984597\) & \(0.927913\) & \(3.7103\cdot10^{-3}\) & \(0.579440\) & \(1.295668\) \\
\(76\)   & \(0.994117\) & \(0.955392\) & \(8.4485\cdot10^{-4}\) & \(0.559760\) & \(1.251662\) \\
\(199\)  & \(0.997753\) & \(0.972463\) & \(1.9553\cdot10^{-4}\) & \(0.548906\) & \(1.227392\) \\
\(521\)  & \(0.999142\) & \(0.983009\) & \(4.5637\cdot10^{-5}\) & \(0.542716\) & \(1.213550\) \\
\(1364\) & \(0.999672\) & \(0.989514\) & \(1.0701\cdot10^{-5}\) & \(0.539094\) & \(1.205451\) \\
\(3571\) & \(0.999875\) & \(0.993526\) & \(2.5162\cdot10^{-6}\) & \(0.536935\) & \(1.200624\) \\
\(9349\) & \(0.999952\) & \(0.996002\) & \(5.9254\cdot10^{-7}\) & \(0.535632\) & \(1.197709\) \\
\midrule
conjectured limit & \(1\) & \(1\) & -- & \(0.533573\) & \(1.193107\) \\
\bottomrule
\end{tabularx}
\end{center}

The minimum of \(n^{3/2}(\W(u)-1)\) over \(4\le n\le10^4\) is
\(0.5356319\), at \(n=9349\); the column decreases monotonically towards
the value conjectured in (34). Every entry exceeds the proved bound
\(\frac{\sqrt2}{25}=0.0565685\) of (18), by a factor between \(12\) and
\(9.5\).

For ordinary \(n\), law (27) is already accurate at moderate size:

\begin{center}
\small
\renewcommand{\arraystretch}{1.3}
\begin{tabularx}{0.86\textwidth}{@{}rXXX@{}}
\toprule
\(n\) & \(\{np\}\) & \(n\bigl(\W(u)-1\bigr)\) & \(\Phi(\{np\})\) \\
\midrule
\(37\)   & \(0.773452\) & \(0.4501191\) & \(0.4417257\) \\
\(100\)  & \(0.360680\) & \(0.6622113\) & \(0.6536000\) \\
\(250\)  & \(0.901699\) & \(0.2845883\) & \(0.2839127\) \\
\(777\)  & \(0.242482\) & \(0.6331557\) & \(0.6319304\) \\
\(2000\) & \(0.213596\) & \(0.6180249\) & \(0.6175447\) \\
\(5000\) & \(0.033989\) & \(0.3316817\) & \(0.3315599\) \\
\bottomrule
\end{tabularx}
\end{center}

\textbf{Status of the numerical evidence.} For \(n=4,8,12,20,30,50\) a
general constrained minimisation of \(\W\) over all \(n\)-good arrays
(SLSQP with auxiliary bounds \(m\le u_i\le M\), many random starts)
returns the value from (20) to seven or eight digits, and no narrower
configuration was found for \(4\le n\le30\). This is local optimisation
from random starts, so it is evidence for the conjecture and not a
certificate: it rules out nothing globally. Separately, (19) is narrower
than the Section 4 patch by a factor of order \(n\), for instance
\(7.48\cdot10^{-7}\) against \(2.22\cdot10^{-3}\) at \(n=15127\), which
shows only that the patch is not extremal.

\section*{References}

\begin{enumerate}[label={[\arabic*]}]
\item A.~Hurwitz, \emph{\"Uber die angen\"aherte Darstellung der
Irrationalzahlen durch rationale Br\"uche}, Math.\ Ann.\ \textbf{39}
(1891), 279--284. (Used for (33); the special case needed here also
follows from the exact identity (31).)
\item A.~Ya.~Khinchin, \emph{Continued Fractions}, University of Chicago
Press, 1964. (Convergents, equidistribution of \(\{np\}\), Section 4.)
\item M.~G.~Krein and A.~A.~Nudelman, \emph{The Markov Moment Problem and
Extremal Problems}, Amer.\ Math.\ Soc., 1977. (The moment-space
framework that a proof of (27) would need.)
\end{enumerate}

\end{document}
